% FONTE TEMA https://github.com/matze/mtheme
\documentclass[aspectratio=1610]{beamer}
%\documentclass[aspectratio=1610, handout]{beamer}
\usepackage[utf8]{inputenc}
\usepackage{ragged2e}
\usepackage{xcolor}
\usepackage[italian]{babel}
\usepackage{multirow}
\usepackage{silence}
\usepackage{listings}
\definecolor{verde}{rgb}{0,0.6,0}
\definecolor{grigio}{rgb}{0.5,0.5,0.5}
\definecolor{rosa}{rgb}{0.58,0,0.82}
\definecolor{backgroud}{rgb}{0.95,0.95,0.92}
\lstdefinestyle{stilepython}{
    backgroundcolor=\color{backgroud},
    commentstyle=\color{verde},
    keywordstyle=\color{magenta},
    numberstyle=\tiny\color{grigio},
    stringstyle=\color{rosa},
    basicstyle=\ttfamily\footnotesize,
    breakatwhitespace=false,
    numbers=left,
    numbersep=5pt,
    showspaces=false,
    showstringspaces=false,
    showtabs=false,
    literate=
    {à}{{\'a}}1
    {è}{{\'e}}1
    {ì}{{\'i}}1
    {ò}{{\'o}}1
    {ù}{{\'u}}1
}
\lstset{style=stilepython}

\WarningFilter{beamer}{}
\WarningFilter{metropolis}{}
\usetheme[progressbar=frametitle,titleformat=smallcaps]{metropolis}
\setbeamertemplate{frame numbering}[fraction]
\setbeamercovered{dynamic}
\definecolor{rosso}{RGB}{255, 0, 0}
\definecolor{giallo}{RGB}{254,212,23}
\hypersetup{colorlinks=true,linkcolor=black,urlcolor=rosso}
\setbeamercolor{palette primary}{fg=black, bg=giallo}
\setbeamercolor{background canvas}{bg=white}
\setbeamercolor{normal text}{fg=black}
\setbeamercolor{progress bar}{fg=rosso}
\setbeamercolor{framesubtitle}{fg=rosso}
\setbeamercolor{normal text .dimmed}{fg=giallo}
\setbeamercolor{block title alerted}{fg=rosso, bg=giallo}
\setbeamerfont{caption}{size=\tiny}
\setbeamerfont{caption name}{size=\tiny}
\setlength{\abovecaptionskip}{0pt}
\makeatletter
\metroset{block=fill}
\setlength{\metropolis@progressinheadfoot@linewidth}{1pt} 
\setlength{\metropolis@progressonsectionpage@linewidth}{1pt}
\setlength{\metropolis@titleseparator@linewidth}{1pt}
\makeatother

\title{FUNZIONI, METODI, MODULI E LIBRERIE}
\subtitle{Elementi fondamentali Python}
\date{}
\institute{\textit{
        Fonti:
        \begin{itemize}
            \item[-] \href{https://docs.python.org/it/3.12/library/functions.html}{Funzioni Incorporate - Documentazione Ufficiale Python}
            \item[-] \href{https://docs.python.org/it/3.14/builtins/stdtypes.html\#text-and-binary-sequence-type-methods-summary}{Metodi Stringe - Documentazione Ufficiale Python}
            \item[-] \href{https://docs.python.org/it/3.14/library/math.html}{Modulo Math - Documentazione Ufficiale Python}
        \end{itemize}
    }
}

\begin{document}

\begin{frame}[plain, noframenumbering]
    \titlepage
\end{frame}

\section{FUNZIONI}

\begin{frame}{FUNZIONI BUILT-IN}
    \begin{alertblock}{DEFINIZIONE}
        \begin{minipage}{0.98\linewidth}
            \justifying
            Le \textbf{funzioni built-in} sono funzioni predefinite e pronte all'uso, 
            sempre disponibili. Eseguono le operazioni fondamentali ed essenziali del 
            linguaggio (ad esempio: input() e print()). Evitano quindi di dover riscrivere 
            manualmente il codice per le attività di base più frequenti.\\
            \bigskip
            \tiny{\textbf{DOCUMENTAZIONE}}\\
            \tiny{\href{https://docs.python.org/it/3.12/library/functions.html}{Tutte le funzioni Incorporate}}
        \end{minipage}
    \end{alertblock}
\end{frame}

\begin{frame}[fragile]{FUNZIONI}
    \begin{alertblock}{DEFINIZIONE}
        \begin{minipage}{0.98\linewidth}
            \justifying
            \textbf{Una funzione} in Python è un \textbf{blocco di codice riutilizzabile} 
            che \textbf{viene eseguito solo quando viene invocato}. Si \textbf{definisce} mediante la parola chiave \textbf{def}. 
            Può accettare dati in ingresso (parametri o argomenti) e restituire un risultato tramite \textbf{return}. 
            Serve per strutturare il programma in modo modulare, evitando duplicazioni di codice e migliorandone la leggibilità.
\begin{lstlisting}[language=Python]
 def area_rettangolo(base, altezza):
    """Calcola e restituisce l'area di un rettangolo"""
    area = base * altezza
    return area
 
 # Chiamata alla funzione
 risultato = area_rettangolo(5,3)
 print(risultato) # output: 15
\end{lstlisting}
        \end{minipage}
    \end{alertblock}
\end{frame}

\section{METODI}

\begin{frame}[fragile]{METODO}
    \begin{alertblock}{DEFINIZIONE}
        \begin{minipage}{0.98\linewidth}
            \justifying
            Un \textbf{metodo} in Python è una funzione definita 
            all'\textbf{interno di una classe} e \textbf{associata a uno specifico 
            oggetto}. Si \textbf{richiama} usando la \textbf{notazione puntata} ovvero:\\ 
            \bigskip
            \textbf{oggetto.metodo()}\\
            \bigskip
            Un metodo può leggere e modificare lo stato e gli attributi 
            dell'oggetto a cui appartiene.\\
            \bigskip
\begin{lstlisting}[language=Python]
 txt = "python"
 txt = txt.upper()
 print(txt) # output: PYTHON
\end{lstlisting}
        \end{minipage}
    \end{alertblock}
\end{frame}

\begin{frame}{METODI STR}
    \begin{alertblock}{DEFINIZIONE}
        \begin{minipage}{0.98\linewidth}
            \justifying
            Nel tipo \textbf{str} di Python, un metodo è una funzione integrata associata 
            a un oggetto stringa, richiamabile tramite la notazione con il punto. Serve per 
            analizzare, formattare o trasformare il testo. \textbf{Poiché le stringhe sono immutabili, 
            un metodo str non altera mai l'oggetto originale, ma restituisce sempre un nuovo valore}.\\
            \bigskip
            \tiny{\textbf{DOCUMENTAZIONE}}\\
            \tiny{\href{https://docs.python.org/it/3.14/builtins/stdtypes.html\#text-and-binary-sequence-type-methods-summary}{Tutti i metodi di str}}
        \end{minipage}
    \end{alertblock}
\end{frame}

\section{MODULI E LIBRERIE}

\begin{frame}{MODULI E LIBRERIE }
    \begin{alertblock}{DEFINIZIONI}
        \begin{minipage}{0.98\linewidth}
            \justifying
            Un \textbf{modulo} in Python è un \textbf{file contenente codice} (funzioni, classi e variabili) 
            salvato con estensione \textbf{.py}. Il suo scopo principale è organizzare il 
            codice e favorirne il riutilizzo, dividendo un programma in sezioni logiche e 
            gestibili. \textbf{Per accedere ai contenuti di un modulo di un altro script si 
            utilizza la parola chiave: import} (ad esempio, import math).\\
            \pause
            Una \textbf{libreria} in Python è una \textbf{raccolta più ampia e complessa 
            di pacchetti e moduli correlati}, progettata per offrire funzionalità complete su 
            un determinato argomento. Mentre un modulo è un singolo file .py, una libreria 
            unisce molti moduli e risorse pronte all'uso. Viene solitamente installata tramite un gestore (es: \textbf{pip}) 
            e poi importata nei propri progetti.\\
            \bigskip
            \tiny{\textbf{DOCUMENTAZIONE}}\\
            \tiny{\href{https://docs.python.org/it/3.14/library/math.html}{Modulo math}}
        \end{minipage}
    \end{alertblock}
\end{frame}

\begin{frame}[fragile]{IMPORTARE MODULI/FUNZIONI}
    \begin{alertblock}{IMPORT/FROM}
        \begin{minipage}{0.98\linewidth}
            \justifying
            Se viene importato tutto il modulo, i metodi dei moduli devono essere richiamati 
            tramite la notazione puntata (che ha come radice il nome del modulo)
\begin{lstlisting}[language=Python]
 import math
 print(math.sqrt(2))
\end{lstlisting}
            Per \textbf{importare un singolo metodo} e poterlo utilizzare come se fosse una 
            funzione built-in è necessario richiamarlo tramite la sintassi: \textbf{from}
            modulo \textbf{import} metodo.
\begin{lstlisting}[language=Python]
 from math import sqrt
 print(sqrt(2))
\end{lstlisting}
        \end{minipage}
    \end{alertblock}
\end{frame}

\end{document}